\scanpage{125}
\mspage{112}
\diagram{\((6.3.6)\)\quad
\begin{tikzcd}[ampersand replacement=\&,column sep=small,row sep=small]
 A'_{I'}\arrow[rr,"{\alpha_{I'}}"]\arrow[dr,"{i'_{I'}}"']\arrow[dd,"{\lambda_A}"'] \& \& X'_{I'}\arrow[dr,"{p_{I'}}"]\arrow[dd,"{\lambda_X}"{pos=.8}] \& \\
 \& B'_{I'}\arrow[rr,"{\beta_{I'}}",crossing over]\arrow[ur,dashed,"{h'_{I'}}"]\arrow[dd,"{\lambda_B}"{pos=.75},crossing over] \& \& Y'_{I'}\arrow[dd,dash,"{\lambda_Y}"]\\
 A_0\arrow[rr,"{\deleted{\unsure{\alpha}}}"{pos=.2},"{\alpha_0}"{pos=.75}]\arrow[dr,"{i_0}"'] \& \& X_0\arrow[dr,"{p_0}"] \& \\
 \& B_0\arrow[rr,"{\beta_0}"]\arrow[ur,dotted,"{h_0}"] \& \& Y_0
\end{tikzcd}}

\margin{NB. La condition $p\in\Psi\langle I\rangle$ s'exprime par $p'\in\Psi(I')$ et $\boxed{p_0\in\Psi}$}
D'autre part, la question de prolonger
$(I',h')\in\unsure{\mathcal H|H}$ en $(I,h)$, revient à celle
de trouver $h_0\colon B_0\longrightarrow X_0$ dans $M$,
de façon qu'avec $h'_{I'}$ \unsure{cel|cela} fasse encore
un diagramme commutatif \unsure{de|dans} (6.3.6),
ce qui revient à dire que \inserted{d'une part} les deux
triangles formés dans le carré du
bas sont commutatifs :
\[
 p_0 h_0=\beta_0,\qquad h_0 i_0=\alpha_0,\tag*{(6.3.7)}
\]
et d'autre part le carré diagonal défini
par $h'_{I'}$, $h_0$ soit commutatif (ce qui
exprime qu'on a une flèche $h$ de
$B$ dans $X$) :
\[
 \lambda_X h'_{I'}=h_0\lambda_B.\tag*{(6.3.8)}
\]
% archive sheet number: 124

